\begingroup


\fontsize{12pt}{10pt}
\selectfont
\renewcommand{\arraystretch}{1}



\phantomsection
\pdfbookmark[1]{Задание на ВКР}{Задание на ВКР}

\vkrStandartHead

\vspace{1cm}

\hfill%
\begin{minipage}{0.5\linewidth}
	\raggedleft
	\textbf{<<УТВЕРЖДАЮ>>}\\
	Заведующий кафедрой ИТиВС\\
	\vkrKafedraHeadTitles\space\vkrKafedraHeadShort\\[3ex]
	\underline{\hspace{30mm}}\\
	<<\underline{\hspace{7mm}}>>\underline{\hspace{30mm}} 2026 г.
\end{minipage}

\vspace{0.5cm}

\begin{center}
	{\fontsize{14pt}{14pt}\selectfont\textbf{Задание}}

	на выполнение выпускной квалификационной работы \\
	по направлению подготовки \\
	09.04.04 <<Программная инженерия>>\\
	Профиль: <<Технологии разработки интеллектуальных систем и программных комплексов>>
\end{center}

{
	\centering\fontsize{14pt}{12pt}\selectfont
	\begin{tabular} {l l l}
		Студент группы \vkrStudentGroup    	& &  	Научный руководитель \\
		\vkrStudentShort					& & 	\insertPersonInfo{Advisor} \vkrAdvisorShort\\
	\end{tabular}\\
}

\begin{center}
	\bfseries\fontsize{14pt}{12pt}\selectfont
	\parbox{0.9\linewidth}{\centering Тема: <<\vkrTheme}

\end{center}

\vfill

\noindent Тема утверждена приказом <<\underline{\hspace{7mm}}>>.\underline{\hspace{30mm}} 202\_ г. № \_\_\_\_

\noindent Срок сдачи ВКР на кафедру <<\underline{\hspace{7mm}}>>.\underline{\hspace{30mm}} 202\_ г.

\newpage
\fontsize{12pt}{16pt}
\selectfont
\begin{enumerate} [nosep,font=\bfseries, left=0pt]
	\item \textbf{Описание задания на выполнение ВКР }
		\begin{enumerate} [nosep, label*=\arabic*.,font=\bfseries]
			\item \textbf{Тип ВКР --- исследовательская работа. }
			\item \textbf{Цель исследования --- повышение эффективности деятельности учебного заведения. }
			\item \textbf{Объект исследования --- программная поддержка процесса создания расписания учебных занятий. }
			\item \textbf{Предмет исследования --- архитектура и функционирование клиент-серверного веб-приложения для создания расписания учебных занятий с учётом личных предпочтений преподавательского состава. }
			\item \textbf{Методы исследования ---  системный анализ, процессно-ориентированный подход и функциональное моделирование процессов, объектно-ориентированный анализ. }
			\item \textbf{Задачи исследования:}
				\begin{enumerate} [nosep, label*=\arabic*.]
					\item Провести исследование различных методов автоматического составления расписаний учебных занятий. Сделать обоснованные выводы о наличии необходимости в разработке своего метода. Провести анализ правил составления расписания МГТУ «СТАНКИН».
					\item Обосновать требования для клиент-серверного веб-приложения для создания расписания учебных занятий.
					\item Разработать проект клиент-серверного веб-приложения для создания расписания учебных занятий.
					\item Разработать клиент-серверное веб-приложение для создания расписания учебных занятий в соответствии с разработанным проектом.
				\end{enumerate}
		\end{enumerate}
	\item \textbf{Требования к выполнению ВКР}
		\begin{enumerate} [nosep, label*=\arabic*.,font=\bfseries]

			\item \textbf{Соблюдение требований законодательной базы и стандартов}

				\begin{enumerate} [nosep, label*=\arabic*.]
					%===============================
					%  Выбрать  свою программу
					%===============================

					\item Образовательная программа высшего образования в магистратуре ФГБОУ ВО «МГТУ «СТАНКИН» по направлению подготовки 09.04.04 «Программная инженерия» для направленности (профиля) подготовки «Технологии разработки интеллектуальных систем и программных комплексов» (утв. 31.05.2021).

					\if 0  %это такой хитрый способ сделать многострочный коментарий

					\item  Образовательная программа высшего образования по направлению подготовки 09.03.01 Информатика и вычислительная техника : направленность (профиль) «Разработка программных комплексов в рамках цифровой трансформации деятельности предприятий» : уровень высшего образования бакалавриат : форма обучения очная  / ФГБОУ ВО «МГТУ «СТАНКИН». --- Москва, 2022 --- одобрена Учёным советом Университета 23.04.2022, утв. ректором Серебренным В.В 29.04.2022  --- (изменена и дополнена).

					\item Образовательная программа высшего образования по направлению подготовки 09.03.04 Программная инженерия : направленность (профиль) «Системный анализ и проектирование программных комплексов» : уровень высшего образования бакалавриат : форма обучения очная / ФГБОУ ВО «МГТУ «СТАНКИН». --- Москва, 2022 --- одобрена Учёным советом Университета 23.04.2022, утв. ректором Серебренным В.В 29.04.2022  --- (изменена и дополнена).

					\item Образовательная программа высшего образования по направлению подготовки 09.04.01 Информатика и вычислительная техника : направленность (профиль) «Управление программными продуктами и проектами» : уровень высшего образования магистратура : форма обучения очная / ФГБОУ ВО «МГТУ «СТАНКИН». --- Москва, 2022 --- одобрена Учёным советом Университета 23.04.2022, утв. ректором Серебренным В.В 29.04.2022  --- (изменена и дополнена).

					\item Образовательная программа высшего образования по направлению подготовки 09.04.03  Программная инженерия : направленность (профиль) «Технологии разработки интеллектуальных систем и программных комплексов» : уровень высшего образования магистратура : форма обучения очная / ФГБОУ ВО «МГТУ «СТАНКИН». --- Москва, 2022 --- одобрена Учёным советом Университета 23.04.2022, утв. ректором Серебренным В.В 29.04.2022  --- (изменена и дополнена).

					\fi

					%===============================
					%  Общая часть
					%===============================

					\item   П 01-04/7/2024. Положение о государственной итоговой аттестации по образовательным программам высшего образования --- программам бакалавриата, программам специалитета и программам магистратуры  --- утверждено приказом ректора от 03.04.2024\,г. №700/1. ---  ФГБОУ ВО «МГТУ «СТАНКИН».

					\item П 01-04/9/2024. Положение о выпускной квалификационной работе обучающихся по образовательным программам высшего образования --- программам бакалавриата, программам специалитета, программам магистратуры   --- утверждено приказом ректора от 02.09.2024. №700/1. --- ФГБОУ ВО «МГТУ «СТАНКИН».

				\end{enumerate}

			\item \textbf{Дополнительные требования}

				\begin{enumerate} [nosep, label*=\arabic*.]
					\item 	Оформление раздела «Список литературы» по национальному стандарту ГОСТ Р 7.0.100-2018 <<Библиографическая запись. Библиографическое описание. Общие требования и правила составления>>
					\item Результаты исследования должны быть опубликованы в виде научных статей и тезисов докладов (не менее 2), а также доложены на научно-технических конференциях и семинарах.
				\end{enumerate}
		\end{enumerate}

	\item \textbf{Срок сдачи оформленной квалификационной работы на кафедру --- вторая декада мая 2026\,г. }

\end{enumerate}

\vspace{3ex}%


\noindent%
\begin{tabularx}{\linewidth}{ @{} X @{}  l  @{}    }

	Исполнитель 							&  \vkrStudentShort   \\[3ex]

	Научный руководитель             		& 					   \\
    \insertPersonInfoB{Advisor} 	 		&   \vkrAdvisorShort
\end{tabularx}


\endgroup

